\chapter{Supersymmetric Localization and Exact Observables}
\label{ch:volume-iii-supersymmetric-localization-exact-observables}

\section{The Localization Principle}

Let \(M\) be a compact Euclidean manifold and let the theory admit a
fermionic symmetry \(Q\) acting on the fields \(\Phi\).  Assume
\[
  Q^2=\mathcal L_v+\delta_{\Lambda}+\delta_R+\delta_F,
  \label{eq:localization-q-square}
\]
where \(\mathcal L_v\) is a Lie derivative along a vector field \(v\),
\(\delta_{\Lambda}\) is a gauge transformation, and \(\delta_R,\delta_F\) are
R-symmetry and flavor-symmetry transformations.  Let \(\mathcal O\) be an
observable invariant under \(Q\) and under the bosonic transformations on the
right side of \eqref{eq:localization-q-square}.

Choose a fermionic functional \(V[\Phi]\) such that \(QV\) is invariant under
the same bosonic transformations.  Define the deformed expectation value
\[
  \langle\mathcal O\rangle_t
  =
  \frac1{Z_t}\int\mathcal D\Phi\,
  \mathcal O[\Phi]\,
  \exp\left[-S[\Phi]-t\,QV[\Phi]\right].
\]
If the measure is \(Q\)-invariant and integration by parts in field space has
no boundary contribution, then
\[
  \frac{\dd}{\dd t}\langle\mathcal O\rangle_t=0.
\]
Thus the expectation value may be computed at any convenient value of \(t\).
When the bosonic part of \(QV\) is nonnegative, the limit \(t\to\infty\)
localizes the integral to the locus
\begin{equation}
  (QV)_{\rm bos}=0.
  \label{eq:localization-bps-locus}
\end{equation}
This locus is the BPS locus for the chosen supercharge and deformation.

\begin{figure}[H]
  \centering
  \resizebox{0.92\textwidth}{!}{%
  \begin{tikzpicture}[x=1cm,y=1cm,every node/.style={font=\scriptsize}]
    \tikzset{
      arrow/.style={-{Latex[length=2mm]},line width=0.85pt},
      field/.style={draw,rounded corners=6pt,fill=blue!4,line width=0.8pt},
      locus/.style={green!45!black,line width=1.0pt}
    }
    \node[field,minimum width=4.0cm,minimum height=2.4cm] (space) at (-3.0,0)
      {field space};
    \draw[locus] (-4.0,-0.35) .. controls (-3.35,0.45) and (-2.65,0.45) .. (-2.0,-0.35);
    \node[green!45!black] at (-3.0,-0.80) {\((QV)_{\rm bos}=0\)};
    \foreach \x/\y in {-4.20/0.65,-3.55/0.95,-2.35/0.78,-1.85/0.32} {
      \draw[orange!85!black,-{Latex[length=1.8mm]},line width=0.75pt]
        (\x,\y)--({0.82*\x-0.54},{0.55*\y-0.05});
    }
    \draw[arrow] (-0.55,0)--(0.90,0);
    \node[align=center] at (0.18,0.55) {\(t\to\infty\)};
    \node[draw,rounded corners=4pt,fill=green!8,minimum width=3.4cm,
      minimum height=1.0cm,align=center] at (3.05,0)
      {integral over\\BPS locus};
    \node[blue!65!black,align=center] at (0,-1.45)
      {normal fluctuations give one-loop determinants};
  \end{tikzpicture}
  }
  \caption{A \(Q\)-exact deformation concentrates the path integral on the
  locus where the bosonic part of \(QV\) vanishes.}
  \label{fig:localization-flow-to-bps-locus}
\end{figure}

\section{One-Loop Determinants}

Let the BPS locus have connected components \(\mathcal M_\alpha\).  Near a
point \(a\in\mathcal M_\alpha\), decompose fields into tangent directions
along \(\mathcal M_\alpha\) and normal fluctuations.  The localized path
integral has the structural form
\begin{equation}
  \langle\mathcal O\rangle
  =
  \frac1Z
  \sum_\alpha
  \int_{\mathcal M_\alpha}\dd\mu_\alpha(a)\,
  \mathcal O_{\rm loc}(a)\,
  \ee^{-S_{\rm cl}(a)}
  Z_{\rm 1-loop}(a)
  Z_{\rm nonpert}(a).
  \label{eq:localized-path-integral-general}
\end{equation}
Here \(\dd\mu_\alpha\) is the induced measure on the BPS moduli,
\(\mathcal O_{\rm loc}\) is the restriction of the observable,
\(S_{\rm cl}\) is the original action evaluated on the locus,
\(Z_{\rm 1-loop}\) is the ratio of Gaussian determinants of normal
fluctuations, and \(Z_{\rm nonpert}\) includes pointlike or localized
nonperturbative sectors when they are present.

The determinant is computed from the quadratic operator obtained by expanding
\(QV\) in the normal directions.  Supersymmetry pairs many bosonic and
fermionic eigenmodes.  The unpaired modes are organized by an equivariant
index for the complex generated by \(Q\) and \(Q^2\).  Thus
\[
  Z_{\rm 1-loop}
  =
  \left(
    \frac{\det_{\rm fermion} \mathcal K_F}
         {\det_{\rm boson} \mathcal K_B}
  \right)^{1/2}
\]
after regularization, with the determinant ratio interpreted through the
equivariant spectrum.

\section{Sphere Partition Functions}

For a superconformal theory, placing the theory on a round sphere \(S^D\)
requires coupling to background supergravity fields that preserve a chosen
set of supercharges.  The sphere partition function
\[
  Z_{S^D}(\lambda,\bar\lambda)
\]
depends on exactly marginal couplings and on supersymmetric background
parameters.  Derivatives with respect to such parameters insert integrated
operators in protected multiplets.

In four-dimensional \(\mathcal N=2\) superconformal theories, the
\(S^4\) partition function determines the Kahler potential on the conformal
manifold up to the ordinary Kahler ambiguity:
\begin{equation}
  Z_{S^4}(\lambda,\bar\lambda)
  =
  \exp\left(c_{\rm sph}K(\lambda,\bar\lambda)\right)
  \label{eq:s4-partition-kahler-potential}
\end{equation}
in a convention where \(c_{\rm sph}\) is a fixed numerical normalization.
Changing local supersymmetric counterterms shifts
\[
  K(\lambda,\bar\lambda)
  \mapsto
  K(\lambda,\bar\lambda)+f(\lambda)+\bar f(\bar\lambda),
\]
which is the expected Kahler transformation.  The metric
\[
  G_{a\bar b}=\partial_a\partial_{\bar b}K
\]
is invariant under this ambiguity and agrees with the separated-point
two-point function of exactly marginal operators after the normalization of
operators has been fixed.

\begin{figure}[H]
  \centering
  \resizebox{0.92\textwidth}{!}{%
  \begin{tikzpicture}[x=1cm,y=1cm,every node/.style={font=\scriptsize}]
    \tikzset{
      arrow/.style={-{Latex[length=2mm]},line width=0.85pt},
      box/.style={draw,rounded corners=4pt,line width=0.8pt,fill=blue!4,
        minimum width=3.2cm,minimum height=0.9cm,align=center}
    }
    \node[box] (sphere) at (-4.1,0.9) {supersymmetric\\sphere background};
    \node[box] (loc) at (0,0.9) {localized\\matrix integral};
    \node[box] (z) at (4.1,0.9) {\(Z_{S^D}\)};
    \node[box] (deriv) at (0,-0.9) {coupling\\derivatives};
    \node[box] (metric) at (4.1,-0.9) {conformal-manifold\\metric};
    \draw[arrow] (sphere)--(loc);
    \draw[arrow] (loc)--(z);
    \draw[arrow] (z)--(deriv);
    \draw[arrow] (deriv)--(metric);
    \node[green!45!black,align=center] at (-3.95,-0.95)
      {exact path integral data\\becomes exact CFT data};
  \end{tikzpicture}
  }
  \caption{Supersymmetric sphere partition functions convert localized path
  integrals into exact observables on conformal manifolds.}
  \label{fig:sphere-partition-function-cft-data}
\end{figure}

\section{Wilson Loops from Matrix Models}

In four-dimensional gauge theories with enough supersymmetry, a circular
Wilson loop can preserve the same supercharge used for localization.  For a
connection \(A_\mu\) and scalar fields \(\phi^I\), a typical supersymmetric
loop has the form
\[
  W_R(C)
  =
  \operatorname{Tr}_R
  P\exp\oint_C
  \left(
    \ii A_\mu\dot x^\mu
    +
    |\dot x|\,n_I\phi^I
  \right)\dd s,
\]
where \(n_I\) is a constant unit vector in the scalar R-symmetry space and
\(R\) is a representation of the gauge group.  Localization maps its
expectation value to an insertion in the same finite-dimensional integral
that computes \(Z_{S^4}\).

For \(\mathcal N=4\) super-Yang--Mills with gauge group \(U(N)\), the
localized integral is Gaussian:
\begin{equation}
  Z_{\mathcal N=4}
  =
  \int_{\mathfrak u(N)}\dd a\,
  \exp\left[-\frac{8\pi^2N}{\lambda}\operatorname{Tr}(a^2)\right],
  \label{eq:n4-gaussian-matrix-model}
\end{equation}
with the Vandermonde determinant included in \(\dd a\) after diagonalization.
For the fundamental circular loop,
\[
  W_\square(a)=\frac1N\operatorname{Tr}\ee^{2\pi a}.
\]
At large \(N\) with fixed 't Hooft coupling \(\lambda\), the expectation value
is
\begin{equation}
  \langle W_\square\rangle
  =
  \frac{2}{\sqrt\lambda}I_1(\sqrt\lambda),
  \label{eq:n4-large-n-wilson-loop}
\end{equation}
where \(I_1\) is the modified Bessel function.  This exact interpolation is a
protected observable of the conformal theory and, at strong coupling, matches
the semiclassical string worldsheet in the holographic description.

\begin{figure}[H]
  \centering
  \resizebox{0.90\textwidth}{!}{%
  \begin{tikzpicture}[x=1cm,y=1cm,every node/.style={font=\scriptsize}]
    \tikzset{
      arrow/.style={-{Latex[length=2mm]},line width=0.85pt},
      eig/.style={circle,fill=blue!65!black,inner sep=1.7pt},
      loop/.style={violet!80!black,line width=1.0pt}
    }
    \begin{scope}[shift={(-3.8,0)}]
      \draw[loop] (0,0) circle (0.86);
      \node[violet!80!black] at (0,1.22) {BPS circle};
      \node at (0,-1.22) {\(W_\square(C)\)};
    \end{scope}
    \draw[arrow] (-2.35,0)--(-1.05,0);
    \node[align=center] at (-1.70,0.48) {localize};
    \begin{scope}[shift={(0.9,0)}]
      \draw[->,thick] (-1.35,-1.0)--(1.45,-1.0) node[right] {\(a\)};
      \draw[blue!65!black,line width=1.0pt,domain=-1:1,samples=80]
        plot (\x,{0.88*sqrt(1-\x*\x)-1.0});
      \foreach \x in {-0.85,-0.45,-0.15,0.25,0.62,0.92} {
        \node[eig] at (\x,{-1.0+0.88*sqrt(1-\x*\x)}) {};
      }
      \node[blue!65!black] at (0,0.33) {eigenvalue density};
      \node at (0,-1.42) {Gaussian matrix model};
    \end{scope}
    \draw[arrow] (2.65,0)--(3.75,0);
    \node[align=center,green!45!black] at (4.75,0.25)
      {\(\displaystyle \frac{2I_1(\sqrt\lambda)}{\sqrt\lambda}\)};
  \end{tikzpicture}
  }
  \caption{A supersymmetric circular Wilson loop in \(\mathcal N=4\)
  super-Yang--Mills is computed by a Gaussian matrix model.}
  \label{fig:localization-wilson-loop-matrix-model}
\end{figure}

\section{Assumptions and Output}

Localization is an exact deformation argument.  Its output is reliable when
the chosen contour, boundary behavior in field space, regularization, and
supersymmetric counterterms are specified.  These choices determine the
scheme-dependent parts of the answer, while separated-point CFT data,
protected OPE coefficients, indices, and invariant conformal-manifold metrics
are extracted from scheme-independent combinations.

The method supplies exact data that can be inserted into the conformal
bootstrap, into duality constraints, and into the large-\(N\) holographic
expansion.  The next chapter develops this interface systematically.
